\providecommand{\cZ}{\mathcal{Z}}
\section{Beyond the invertible sector}\label{sec:singular}
This section studies series with singular minimal letter matrices. It opens with a local necessary condition derived from Theorem~\ref{thm:main} (Corollary~\ref{cor:local}). Theorem~\ref{thm:rankone} characterizes CCRA-computability exactly when every singular minimal letter matrix has rank at most one. Theorem~\ref{thm:affinecontext} is a new necessary condition of arithmetic nature, derived from an inverse-weight lemma for compressed machines (Proposition~\ref{prop:invweight}); it implies that CCRA over $\Q$ are not closed under reversal (Corollary~\ref{cor:reversal}), which proves Conjecture~51 of~\cite{MPSfull}, and that the candidate of~\cite[Example~52]{MPSfull} is a witness. A companion lifting statement (Proposition~\ref{prop:lifting}) controls translations as well and yields a second obstruction (Theorem~\ref{thm:unipotentcontext}), which shows that a block-symmetric elementary symmetric function, equal to its own reversal, is not CCRA-computable although all local groups are commutative. Theorem~\ref{thm:onereset} shows that these necessary conditions are also sufficient for series over $\{a,\#\}$ with $M_a$ invertible and all words containing $\#$ of rank two. Proposition~\ref{prop:interact} shows that the normal form~(C2) of Theorem~\ref{thm:main} fails beyond the invertible sector: registers must interact. Throughout, $f$ has minimal representation $(\lambda,(M_a)_a,\gamma)$ of dimension $d$, words act on row vectors from the right, and for a subspace $X\subseteq\Q^{1\times d}$ we write $X^\perp=\{v\in\Q^d:Xv=0\}$.

\subsection{A local necessary condition}\label{ssec:local}
Theorem~\ref{thm:main} constrains singular series locally.
\begin{corollary}[Local necessary condition]\label{cor:local}
Let $f$ be computed by a CCRA, with minimal representation $(\lambda,(M_a)_a,\gamma)$ of dimension $d$. Let $e$ be a word with $\operatorname{rank}M_e^2=\operatorname{rank}M_e$, let $R=\Q^{1\times d}M_e$ be the row space of $M_e$, and let $Z$ be a finite set of words with $\operatorname{rank}(M_eM_zM_e)=\operatorname{rank}M_e$ for all $z\in Z$. Then the maps $\nu_z:R\to R$, $r\mapsto rM_eM_zM_e$, are bijective, and the Zariski closure of the group they generate satisfies condition~(C3) of Theorem~\ref{thm:main}: its identity component is triangularizable over $\Q$ and its unipotent elements commute.
\end{corollary}
\begin{proof}
Write $N_z=M_eM_zM_e$. The rank condition on $e$ gives $\Q^{1\times d}=R\oplus\{r:rM_e=0\}$ and $RM_e=R$, so $RN_z=RM_zM_e=\Q^{1\times d}N_z$ has dimension $\operatorname{rank}M_e=\dim R$; hence $\nu_z$ is a bijection of $R$. Let $\cH\subseteq\GL(R)$ be the closure of the group generated by the $\nu_z$.

For words $u,v$ consider $g_{u,v}(z_1\cdots z_n)=f(ue\cdot ez_1e\cdots ez_ne\cdot ev)=c_v(\rho_uN_{z_1}\cdots N_{z_n})$, where $\rho_u=\lambda M_uM_e\in R$ and $c_v:r\mapsto rM_eM_v\gamma$. Reading $ue$, then the blocks $eze$, and finally $ev$ before the output shows that $g_{u,v}$ is computed by a CCRA over the alphabet $Z$. By the row version of Section~\ref{sec:prelim}, its minimal representation is $X_u/Y_{u,v}$ with $X_u=\spn\{\rho_uN_w:w\in Z^*\}$ and $Y_{u,v}=\{r\in X_u:c_v(rN_w)=0\ \forall w\in Z^*\}$, where $N_{z_1\cdots z_n}=N_{z_1}\cdots N_{z_n}$; its letter operators are induced by the bijections $\nu_z$, so $g_{u,v}$ is invertible. By Theorem~\ref{thm:main}, the closure $\cH_{u,v}$ of the induced group on $X_u/Y_{u,v}$ satisfies~(C3), and restriction maps $\cH$ onto $\cH_{u,v}$ by (A2).

For matrices write $[X,Y]=XY-YX$. If $h_1,h_2\in\cH$ are unipotent, their images in $\cH_{u,v}$ commute, so $[h_1,h_2]$ maps $X_u$ into $Y_{u,v}$; with $w=\eps$ this gives $c_v(\rho_u[h_1,h_2])=0$ for all $u,v$. The rows $\rho_u$ span $R$ (as the $\lambda M_u$ span the dual space), and the functionals $c_v$ separate $R$ (if $rM_eM_v\gamma=0$ for all $v$ then $rM_e=0$, so $r=0$). Hence $[h_1,h_2]=0$. For triangularizability, $\bigcap_vY_{u,v}=0$ by the same separation, so $X_u$ embeds $\cH^\circ$-equivariantly in a finite sum $\bigoplus_vX_u/Y_{u,v}$, each of which carries an $\cH^\circ$-invariant complete rational flag; thus each $X_u$ has one, and so does $R=\sum_uX_u$, a quotient of $\bigoplus_uX_u$.
\end{proof}
For $e=\eps$ the corollary says that any finite set of words whose minimal matrices are invertible generates a group satisfying~(C3), even if other letters are singular. For general $e$ it constrains the action of the words $eze$ on the row space of every word matrix $M_e$ with $\operatorname{rank}M_e^2=\operatorname{rank}M_e$.

\begin{lemma}\label{lem:quotients}
If $f$ is computed by a CCRA, $v,v'\in\Sigma^*$ and $\Sigma'\subseteq\Sigma$, then $u\mapsto f(vuv')$ on $\Sigma'^*$ is computed by a CCRA.
\end{lemma}
\begin{proof}
Start in $\delta(s_0,v)$ with valuation $P_{s_0,v}(x^{(0)})$, keep the transitions on $\Sigma'$, and in state $s$ output $\omega_{\delta(s,v')}\circ P_{s,v'}$, which is copyless as a composite of copyless maps.
\end{proof}

\subsection{Rank-one resets}
\begin{theorem}\label{thm:rankone}
Suppose that every minimal letter matrix of $f$ is invertible or has rank at most one. Let $\Sigma_0$ be the set of letters with invertible matrix and $\cG_0\subseteq\GL_d(\Qb)$ the Zariski closure of the group generated by $\{M_a:a\in\Sigma_0\}$. Then $f$ is CCRA-computable if and only if $\cG_0^\circ$ is triangularizable over $\Q$ and any two unipotent elements of $\cG_0$ commute. Consequently CCRA-computability is decidable in this class, and the CCRA series in this class are closed under reversal.
\end{theorem}
\begin{proof}
Let $\Sigma_1=\Sigma\setminus\Sigma_0$. For $\sigma\in\Sigma_1$ write $M_\sigma=c_\sigma r_\sigma$ with a column $c_\sigma$ and a row $r_\sigma$ (both zero if $M_\sigma=0$), and put $r_\triangleright=\lambda$, $c_\triangleleft=\gamma$. For $i\in\Sigma_1\cup\{\triangleright\}$ and $j\in\Sigma_1\cup\{\triangleleft\}$ define $g_{ij}:\Sigma_0^*\to\Q$, $g_{ij}(u)=r_iM_uc_j$. Every word has a unique factorization $w=u_0\sigma_1u_1\cdots\sigma_ku_k$ with $u_m\in\Sigma_0^*$ and $\sigma_m\in\Sigma_1$, and
\begin{equation}\label{eq:rankonefactor}
f(w)=g_{\triangleright\sigma_1}(u_0)\,g_{\sigma_1\sigma_2}(u_1)\cdots g_{\sigma_k\triangleleft}(u_k).
\end{equation}
The series $g_{ij}$ is computed by $(r_i,(M_a)_{a\in\Sigma_0},c_j)$, so its minimal representation is an invariant subquotient $X_i/Y_{ij}$ of the row space, with $X_i=\spn\{r_iM_u:u\in\Sigma_0^*\}$ and $Y_{ij}=\{x\in X_i:xM_uc_j=0\ \forall u\in\Sigma_0^*\}$ (row version of Section~\ref{sec:prelim}). Its letter operators are induced by invertible matrices, so $g_{ij}$ is invertible, and by (A2) the Zariski closure $\cG_{ij}$ of its group is the image of $\cG_0$ under the induced-action morphism.

\emph{Sufficiency.} Condition~(C3) passes from $\cG_0$ to $\cG_{ij}$: identity components map onto identity components (A2), the traces of a complete $\cG_0^\circ$-invariant rational flag form a chain of invariant rational subspaces of $X_i/Y_{ij}$ with quotients of dimension at most one, and unipotent elements of $\cG_{ij}$ are images of unipotent elements of $\cG_0$ (A2). By Theorem~\ref{thm:main}, each $g_{ij}$ is computed by a CCRA $\mathcal C_{ij}$ over $\Sigma_0$. The machine for $f$ keeps in its control the last letter $i\in\Sigma_1\cup\{\triangleright\}$ read and the control states of the machines $\mathcal C_{ij}$, $j\in\Sigma_1\cup\{\triangleleft\}$; it has an accumulator $\zeta$, initially $1$, and for each $j$ a block $\mathcal B_j$ of registers, large enough for every $\mathcal C_{ij}$, holding the registers of $\mathcal C_{ij}$ for the current $i$. A letter $a\in\Sigma_0$ updates each block $\mathcal B_j$ by $\mathcal C_{ij}$. A letter $\sigma\in\Sigma_1$ performs $\zeta:=\zeta\cdot\omega_{i\sigma}$, where $\omega_{i\sigma}$ is the current output polynomial of $\mathcal C_{i\sigma}$ applied to $\mathcal B_\sigma$, sets $i:=\sigma$, and reinitializes every block $\mathcal B_j$ and its control component to the initial data of $\mathcal C_{\sigma j}$; all other old register values are discarded. In state $i$ the output is $\zeta\cdot\omega_{i\triangleleft}$. Every update uses each old register at most once, and~\eqref{eq:rankonefactor} shows that the machine computes $f$.

\emph{Necessity.} Let $f$ be CCRA-computable. Since the rows $\lambda M_v$ span $\Q^{1\times d}$ and the columns $M_{v'}\gamma$ span $\Q^d$, for nonzero $c_i$ and $r_j$ (otherwise $g_{ij}=0$) there are words $v,v'$ with $\lambda M_vc_i\ne0\ne r_jM_{v'}\gamma$, and $f(v\,i\,u\,j\,v')=(\lambda M_vc_i)\,g_{ij}(u)\,(r_jM_{v'}\gamma)$; similarly for $i=\triangleright$ or $j=\triangleleft$. By Lemma~\ref{lem:quotients} every $g_{ij}$ is CCRA-computable, so $\cG_{ij}$ satisfies~(C3) by Theorem~\ref{thm:main}. Let $h_1,h_2\in\cG_0$ be unipotent. Their images in $\cG_{ij}$ commute, so $[h_1,h_2]$ maps $X_i$ into $\bigcap_jY_{ij}$. If $x\in\bigcap_jY_{ij}$, then $xM_w\gamma=0$ for every word $w$: either $w\in\Sigma_0^*$, or $w=u\sigma w'$ with $u\in\Sigma_0^*$ and $xM_w\gamma=(xM_uc_\sigma)(r_\sigma M_{w'}\gamma)=0$. Hence $x=0$. As every row $\lambda M_w$ lies in $X_\triangleright$ or, if $\sigma$ is the last letter of $w$ in $\Sigma_1$, in $X_\sigma$, we have $\Q^{1\times d}=\sum_iX_i$, and $[h_1,h_2]=0$. For triangularizability, $\cG_0^\circ$ preserves each $X_i$ and each $Y_{ij}$, and its image in $\GL(X_i/Y_{ij})$ is connected, hence contained in $\cG_{ij}^\circ$, which fixes a complete rational flag. Since $\bigcap_jY_{ij}=0$, $X_i$ embeds $\cG_0^\circ$-equivariantly into $\bigoplus_jX_i/Y_{ij}$; intersecting the concatenated flags with the image gives a complete invariant rational flag of $X_i$, and $\Q^{1\times d}=\sum_iX_i$ is a quotient of $\bigoplus_iX_i$. So $\cG_0^\circ$ is triangularizable over $\Q$ on rows, hence on columns.

\emph{Consequences.} Decidability follows from the procedure of Section~\ref{sec:algorithms} applied to $\{M_a:a\in\Sigma_0\}$. The reversal has minimal representation $(\gamma^{\mathsf T},(M_a^{\mathsf T})_a,\lambda^{\mathsf T})$, transposition preserves ranks, and $g\mapsto(g^{\mathsf T})^{-1}$ preserves condition~(C3) as in the proof of Corollary~\ref{cor:consequences}.
\end{proof}

\subsection{Inverse weights and translations}
We return to the setting of Section~\ref{sec:compression}: a monoid $\cY$ acts on $\Q^r$ by copyless maps $P_y$, subspaces $N\subseteq W\subseteq V_r$ are invariant under all $B_y$, $D=W/N$, and every induced operator $\rho(y)$ is invertible. Let $k,L,A,Q$ be as chosen before Lemma~\ref{lem:projection}, and write $e$ for the minimizing element $\eta$ chosen there, let $R_y=Q\circ P_y|_L$ be the affine coordinate permutations of the proof of Proposition~\ref{prop:compression}, and identify $L$ with $\Q^k$ through its free coordinates. Put $\cZ=\{eye:y\in\cY\}$.

\begin{lemma}\label{lem:Zhom}
For $z\in\cZ$ we have $R_z=P_z|_L$, and $R_{zz'}=R_{z'}\circ R_z$ for $z,z'\in\cZ$. Hence $z\mapsto R_z$ is an antihomomorphism from $\cZ$ to the group of affine coordinate permutations of $\Q^k$, and $\cZ_0=\{z\in\cZ:R_z\text{ has trivial permutation part}\}$ is a subsemigroup containing $z^{k!}$ for every $z\in\cZ$.
\end{lemma}
\begin{proof}
$P_{eye}=P_e\circ P_y\circ P_e$ takes values in $L$, where $Q$ is the identity; so $R_z=P_z|_L$ and $P_z(L)\subseteq L$, whence $R_{zz'}=(P_{z'}\circ P_z)|_L=R_{z'}\circ R_z$.
\end{proof}

For $z\in\cZ_0$ write $R_z(t)=(\alpha_i(z)t_i+\beta_i(z))_{i=1}^k$ with $\alpha_i(z)\in\Q^\times$. The multiplier $\alpha_i$ is visible to the machine: if $t\in L$ and $t'=t+\epsilon_i$ differs from $t$ by $1$ in the free coordinate $i$, then
\begin{equation}\label{eq:multiplier}
\alpha_i(z)=P_z(t')_i-P_z(t)_i\qquad(z\in\cZ_0).
\end{equation}

\begin{proposition}[Inverse weights]\label{prop:invweight}
Let $D_0\subset D_1\subset D_2\subseteq D$ be subspaces invariant under $\rho(\cZ_0)$ with $\dim D_1/D_0=\dim D_2/D_1=1$, and let $\kappa_1(z),\kappa_2(z)$ be the scalars by which $\rho(z)$ acts on $D_1/D_0$ and on $D_2/D_1$. Let $\cH$ be the Zariski closure of the group generated by $\rho(\cZ_0)$, and suppose that some unipotent element of $\cH$ induces a non-identity map on $D_2/D_0$. Then there is a free coordinate $i$ with
\[
\alpha_i(z)=\kappa_2(z)/\kappa_1(z)\qquad\text{for all }z\in\cZ_0 .
\]
\end{proposition}
\begin{proof}
Let $\Theta:W''/N''\to D$ be the intertwining isomorphism of Proposition~\ref{prop:compression}, and let $N''\subseteq W_0''\subseteq W_1''\subseteq W_2''\subseteq W''$ be the preimages of $D_0,D_1,D_2$; all are invariant under $B_{R_z}$, $z\in\cZ_0$. Let $\mathcal S$ be the closed subgroup of $\GL(V_k\otimes\Qb)$ stabilizing these five subspaces, $\pi:\mathcal S\to\GL(D)$ and $\pi_2:\mathcal S\to\GL(D_2/D_0)$ the induced-action morphisms, and $\cG_L\subseteq\mathcal S$ the Zariski closure of the group generated by $\{B_{R_z}:z\in\cZ_0\}$. As in Corollary~\ref{cor:onestate}, $\pi(\cG_L)=\cH$. Fix $\bar e_1\in D_1/D_0$ and $\bar e_2\in(D_2/D_0)\setminus(D_1/D_0)$. Every element of $\pi_2(\mathcal S)$ preserves the line $D_1/D_0$, so the unipotent elements of $\pi_2(\cG_L)$ are of the form $u_s:\bar e_1\mapsto\bar e_1$, $\bar e_2\mapsto\bar e_2+s\bar e_1$.

For $b\in\Qb^k$ let $\tau_b(t)=t+b$, and let $U'=\{b:B_{\tau_b}\in\cG_L\}$, a closed subgroup of $\Qb^k$. If $b\in U'$, then $U'\cap\Qb b$ contains $\Z b$, whose Zariski closure is the line $\Qb b$; so $U'$ is a linear subspace. Define $\psi:U'\to\Qb$ by $\pi_2(B_{\tau_b})=u_{\psi(b)}$. It is additive, and $c\mapsto\psi(cb)-c\psi(b)$ is a polynomial function of $c\in\Qb$ vanishing on $\Z$; so $\psi$ is linear.

\emph{$\psi\ne0$.} Let $v\in\cH$ be unipotent and non-identity on $D_2/D_0$. Write $v=\pi(g)$ with $g\in\cG_L$; then $v=v_u=\pi(g_u)$ by (A2), and $g_u\in\cG_L\subseteq\cW_k$ is unipotent, hence $g_u=B_{\tau_b}$ for some $b$ by Lemma~\ref{lem:wk}. Then $b\in U'$, and $\pi_2(g_u)$ is the map induced by $v$ on $D_2/D_0$, so $\psi(b)\ne0$.

\emph{Equivariance.} Let $z\in\cZ_0$ and $\Delta_z=\mathrm{diag}(\alpha_1(z),\ldots,\alpha_k(z))$. A direct computation gives $R_z\circ\tau_b\circ R_z^{-1}=\tau_{\Delta_zb}$, hence $B_{\tau_{\Delta_zb}}=B_{R_z}^{-1}B_{\tau_b}B_{R_z}\in\cG_L$ for $b\in U'$, so $\Delta_zU'=U'$. In the basis $(\bar e_1,\bar e_2)$, $\pi_2(B_{R_z})$ is the upper triangular matrix with diagonal $(\kappa_1(z),\kappa_2(z))$, and for such a matrix $h$ one computes $h^{-1}u_sh=u_{s\kappa_2(z)/\kappa_1(z)}$. Applying $\pi_2$,
\[
\psi(\Delta_zb)=\frac{\kappa_2(z)}{\kappa_1(z)}\,\psi(b)\qquad(b\in U',\ z\in\cZ_0).
\]

\emph{Weight decomposition.} Call free coordinates $i,i'$ equivalent if $\alpha_i=\alpha_{i'}$ on $\cZ_0$. For an equivalence class $C$ let $P_C$ be the coordinate projection onto $\Qb^C$. For each class $C'\ne C$ choose $z_{C'}\in\cZ_0$ with $\alpha_C(z_{C'})\ne\alpha_{C'}(z_{C'})$; then $P_C=\prod_{C'\ne C}(\Delta_{z_{C'}}-\alpha_{C'}(z_{C'}))/(\alpha_C(z_{C'})-\alpha_{C'}(z_{C'}))$ is a polynomial in the operators $\Delta_z$, so $P_CU'\subseteq U'$ and $U'=\bigoplus_CP_CU'$. As $\psi\ne0$, some $b\in P_CU'$ has $\psi(b)\ne0$. Then $\Delta_zb=\alpha_C(z)b$, and the displayed identity gives $\alpha_C(z)=\kappa_2(z)/\kappa_1(z)$ for all $z\in\cZ_0$. Any $i\in C$ is as required.
\end{proof}

In the invertible sector $\kappa_2/\kappa_1$ is a rational character and the conclusion is harmless. In general it says that some register, whose multiplier is computed by the machine through~\eqref{eq:multiplier}, must realize the \emph{inverse} of the weight by which the unipotent part translates. This is the source of the obstruction in Theorem~\ref{thm:affinecontext} below.

The multipliers are only half of the data of the compressed registers; the translations are also rational series, which gives a second constraint. Let $\cG'_L\subseteq\GL(V_k\otimes\Qb)$ be the Zariski closure of the group generated by $\{B_{R_z}:z\in\cZ\}$, let $F'_L=\cG'_L/\cG'^\circ_L$, let $\cZ_1=\{z\in\cZ:B_{R_z}\in\cG'^\circ_L\}$, and let $\cH'$ be the Zariski closure of the group generated by $\rho(\cZ)$. By Lemma~\ref{lem:Zhom}, $z\mapsto B_{R_z}$ is a homomorphism on $\cZ$ ($B_{R_{zz'}}=B_{R_{z'}\circ R_z}=B_{R_z}B_{R_{z'}}$), so $\Psi(z)=B_{R_z}\cG'^\circ_L$ is a homomorphism from $\cZ$ to the finite group $F'_L$, and $\cZ_1=\Psi^{-1}(1)$ is a subsemigroup containing $z^{|F'_L|}$ for every $z\in\cZ$. Since $\cG'^\circ_L$ is connected, it lies in $\cW_k^\circ$, so $\cZ_1\subseteq\cZ_0$.

\begin{proposition}[Rational lifting]\label{prop:lifting}
By the argument of Corollary~\ref{cor:onestate}, the unipotent elements of $\cH'$ commute; let $U_H$ be the unipotent radical of $\cH'^\circ$. There are a maximal torus $T_H$ of $\cH'^\circ$, a linear map $\Lambda:\Qb^k\supseteq U''\to\operatorname{Lie}(U_H)$, and for each free coordinate $i$ a constant $c_i\in\Qb$ such that for every $z\in\cZ_1$ the decomposition $\rho(z)=t(z)u(z)$ with $t(z)\in T_H$, $u(z)\in U_H$ satisfies
\[
\log u(z)=\Lambda(\gamma(z)),\qquad \gamma_i(z)=R_z(c)_i-c_i=\bigl(\alpha_i(z)-1\bigr)c_i+R_z(0)_i ,
\]
and every character of $T_H$ takes on $t(z)$ the value of a Laurent monomial in $\alpha_1(z),\ldots,\alpha_k(z)$. In particular, extending $\Lambda$ linearly to $\Qb^k$, each coordinate of $\log u(z)$ is a $\Qb$-linear combination of the constant $1$ and the rational series $z\mapsto P_z(0)_i$ and $z\mapsto P_z(\epsilon_i)_i-P_z(0)_i$ (computed from fixed valuations in $L$ as in~\eqref{eq:multiplier}).
\end{proposition}
\begin{proof}
Let $\pi:\mathcal S'\to\GL(D)$ be the induced-action morphism on the stabilizer $\mathcal S'$ of $W''$ and $N''$, which contains $\cG'_L$; as in Corollary~\ref{cor:onestate}, $\pi(\cG'_L)=\cH'$ and $\pi(\cG'^\circ_L)=\cH'^\circ$. The group $\cG'^\circ_L$ is connected and solvable, being contained in $\cW_k^\circ$, the group of pullbacks of maps $t\mapsto(a_it_i+b_i)_i$. Choose a maximal torus $T_L$ of $\cG'^\circ_L$. It lies in a maximal torus of $\cW_k^\circ$, and these are the conjugates of the group of pullbacks of diagonal scalings by pullbacks of translations~\cite[10.6]{Bor91}; so there is $c\in\Qb^k$ such that $T_L$ consists of pullbacks of maps $\sigma_a(t)=(a_i(t_i-c_i)+c_i)_i$. The unipotent elements of $\cG'^\circ_L$ form its unipotent radical $U_L$~\cite[10.6]{Bor91}, and by (A3) and Lemma~\ref{lem:wk} they are the pullbacks $B_{\tau_b}$ with $b$ in the subspace $U''=\{b:B_{\tau_b}\in\cG'_L\}$ (a subspace by the argument in the proof of Proposition~\ref{prop:invweight}). For $z\in\cZ_1$ write $R_z=\tau_{\gamma(z)}\circ\sigma_{\alpha(z)}$ with $\gamma(z)=R_z(c)-c$; then $B_{R_z}=B_{\sigma_{\alpha(z)}}B_{\tau_{\gamma(z)}}$. Since $B_{R_z}\in\cG'^\circ_L=T_L\ltimes U_L$ and the decomposition of an element of $\cW_k^\circ$ as a product of a pullback of some $\sigma_a$ and a pullback of a translation is unique, $B_{\sigma_{\alpha(z)}}\in T_L$ and $B_{\tau_{\gamma(z)}}\in U_L$, so $\gamma(z)\in U''$. The image $T_H=\pi(T_L)$ is a maximal torus of $\cH'^\circ$~\cite[11.14]{Bor91}, and $\pi(U_L)=U_H$ by (A2); by uniqueness of the decomposition in $\cH'^\circ=T_H\ltimes U_H$, $t(z)=\pi(B_{\sigma_{\alpha(z)}})$ and $u(z)=\pi(B_{\tau_{\gamma(z)}})$. As $U_H$ is commutative, $\log$ is an isomorphism of $U_H$ onto the vector group $\operatorname{Lie}(U_H)$, and $b\mapsto\log\pi(B_{\tau_b})$ is a homomorphism of vector groups in characteristic zero, hence linear; call it $\Lambda$. Characters of $T_H$ pull back to characters of $T_L$, and these extend to characters of the torus $\{B_{\sigma_a}\}\cong\mathbb G_m^k$ (restriction of characters to a subtorus is onto), which are Laurent monomials in $a$. Finally $R_z(t)_i=\alpha_i(z)t_i+R_z(0)_i$ gives the formula for $\gamma_i(z)$, and $R_z=P_z|_L$ by Lemma~\ref{lem:Zhom}.
\end{proof}

\subsection{Affine contexts and non-closure under reversal}
An \emph{exponential polynomial} is a function $E:\N\to\Qb$ of the form $E(\ell)=\sum_{\mu\in\mathcal M}p_\mu(\ell)\mu^\ell$ with $\mathcal M\subseteq\Qb^\times$ finite and $p_\mu\in\Qb[T]$. We use three standard facts: (E1) the functions $\ell\mapsto\ell^s\mu^\ell$ are linearly independent on every set $\{\ell\ge\ell_0\}$, and the functions $(\ell_1,\ldots,\ell_k)\mapsto\prod_m\ell_m^{s_m}\mu_m^{\ell_m}$ are linearly independent on $\N^k$; (E2) sums and products of exponential polynomials are exponential polynomials; (E3) for a square matrix $Y$ of size $n$, the entries of $Y^\ell$ agree with exponential polynomials for $\ell\ge n$, by the Jordan form. Consequently, if $(\lambda',(M'_a)_a,\gamma')$ represents a rational series $\varphi$ and $w_0,\ldots,w_m,y$ are words, then $\ell\mapsto\varphi(w_0y^\ell w_1y^\ell\cdots y^\ell w_m)$ agrees with an exponential polynomial for large $\ell$.

\begin{lemma}\label{lem:exppoly}
Let $q\in\Qb[T]$ be nonconstant and $E$ an exponential polynomial. Then $q(\ell)E(\ell)=1$ fails for infinitely many $\ell$.
\end{lemma}
\begin{proof}
Otherwise $qE-1=\sum_\mu(qp_\mu)(\ell)\mu^\ell-1$ vanishes for $\ell\ge\ell_0$, and by (E1) the coefficient of $\mu=1$ gives $qp_1=1$ in $\Qb[T]$, which is impossible.
\end{proof}

\begin{definition}\label{def:affinecontext}
An \emph{affine context} for $f$ consists of words $u,x_1,w,x_2$, linearly independent rows $r_1,r_2\in\Q^{1\times d}$ and polynomials $\alpha,\beta\in\Q[T]$ such that, with $b_\ell=x_1w^\ell x_2$:
\begin{enumerate}
\item[(i)] $\lambda M_u\in\Q r_1+\Q r_2$ and $\lambda M_u\notin\Q r_1$;
\item[(ii)] $r_1M_{b_\ell}=r_1$ and $r_2M_{b_\ell}=\beta(\ell)r_1+\alpha(\ell)r_2$ for all $\ell\ge0$;
\item[(iii)] $\alpha$ is nonconstant, and the rational function $\beta/(\alpha-1)$ is nonconstant.
\end{enumerate}
\end{definition}
Thus the blocks $b_\ell$ act on the plane $\Q r_1+\Q r_2$ by $r_1\mapsto r_1$, $r_2\mapsto\beta(\ell)r_1+\alpha(\ell)r_2$: an affine map $s\mapsto s+\beta(\ell)$ on the coefficient of $r_1$, scaled by $\alpha(\ell)$, whose fixed points move with $\ell$. The proof below uses only that the columns $M_v\gamma$ span $\Q^d$, so affine contexts may be exhibited in any representation with this property.

\begin{theorem}\label{thm:affinecontext}
If $f$ has an affine context, then $f$ is not computed by any CCRA over $\Q$.
\end{theorem}
\begin{proof}
Suppose that a CCRA $\mathcal C$ with notation as in Section~\ref{sec:prelim} computes $f$.

\emph{Step 1: control.} Let $n=|S|$ and $p=n!$. For every state $s$ and $j\ge n$ we have $\delta(s,w^{j+p})=\delta(s,w^j)$. Choose $c\ge n$ such that $|\alpha(\ell)|>1$ for all $\ell\ge c$, and let $I=c+p\N$. Then $\tau(s)=\delta(s,b_\ell)$ does not depend on $\ell\in I$, and $\tau^p$ is idempotent. Put $u'=ub_c^p$ and $q=\delta(s_0,u')=\tau^p(\delta(s_0,u))$, so $\tau^p(q)=q$. Let $\cY\subseteq\Sigma^*$ be the monoid generated by the words $b_{\ell_1}\cdots b_{\ell_p}$ with all $\ell_j\in I$. Every $y\in\cY$ fixes $q$, and $\cY$ acts on $\Q^r$ by $P_y=P_{q,y}$.

\emph{Step 2: the local quotient.} Define $W$, $N$ and $K_{u'}$ as in Lemma~\ref{lem:local}, with $\cY$ in place of $L_{S_{\mathrm{ret}}}$ and $u'$ in place of the word $ue$ of that lemma. The proof of Lemma~\ref{lem:local} uses only that the words in question form a monoid fixing $q$, so $W,N$ are invariant under every $B_y$, and $D=W/N\cong\Q^d/K_{u'}$ intertwines $B_y$ with the operator $\rho(y)$ induced by $M_y$. Let $X=\Q r_1+\Q r_2$, which is invariant under each $M_{b_\ell}$ by~(ii). For $s,t\in\Q$ with $t\ne0$, the line $\Q(sr_1+tr_2)$ is $M_{b_\ell}$-invariant iff $s/t=\beta(\ell)/(\alpha(\ell)-1)$, and for $\ell\in I$ the matrix of $M_{b_\ell}$ on $X$ has the distinct eigenvalues $1$ and $\alpha(\ell)$, so $M_{b_\ell}^p$ has the same invariant lines. By~(i), $\lambda M_{u'}=\mu_1r_1+\mu_2r_2$ with $\mu_2\ne0$. If $\spn\{\lambda M_{u'}M_y:y\in\cY\}$ were one-dimensional, the line $\Q\lambda M_{u'}$ would be invariant under every $M_{b_\ell}^p$, $\ell\in I$, contradicting~(iii). Hence this span is $X$, and $K_{u'}=X^\perp$. Choose $v_1,v_2\in\Q^d$ with $r_iv_j=\delta_{ij}$, and let $\bar e_1,\bar e_2$ be their images in $D\cong\Q^d/X^\perp$. The subspace $X^\perp$ is invariant under each $M_{b_\ell}$, and by~(ii) the operator induced by $M_{b_\ell}$ on $\Q^d/X^\perp$ maps $\bar e_1\mapsto\bar e_1+\beta(\ell)\bar e_2$ and $\bar e_2\mapsto\alpha(\ell)\bar e_2$. Hence for $y\in\cY$, $\rho(y)$ fixes the line $\Q\bar e_2$, acts on it by a scalar $\bar\alpha(y)$ equal to the product of the values $\alpha(\ell)$ over the blocks $b_\ell$ of $y$, and acts trivially on $D/\Q\bar e_2$. In particular every $\rho(y)$ is invertible, $\bar\alpha$ is multiplicative, and $|\bar\alpha(y)|>1$ for $y\ne\eps$.

\emph{Step 3: compression.} Proposition~\ref{prop:compression} applies to $\cY$, $W$, $N$. Since $k(P_\eps)=r$ is the largest possible value, we may and do choose the minimizing element (written $e$, the element $\eta$ of Section~\ref{sec:compression}) with $e\ne\eps$. Let $k$, $L$, $R$, $\cZ$, $\cZ_0$ and $\alpha_i$ be as in Lemma~\ref{lem:Zhom}; $k\ge1$ because $\dim D=2$.

\emph{Step 4: a unipotent element.} We claim that $\rho(\cZ_0)$ is not commutative. Otherwise, for $y\in\cY$ the matrices $\rho((eye)^{k!})\in\rho(\cZ_0)$ commute pairwise. The matrix $\rho(eye)=\rho(e)\rho(y)\rho(e)$ has the distinct eigenvalues $1$ and $\bar\alpha(e)^2\bar\alpha(y)$, the latter of absolute value $>1$, so its $1$-eigenline equals that of its $k!$-th power. Commuting diagonalizable $2\times2$ matrices with distinct eigenvalues have the same eigenlines, and here the other eigenline is always $\Q\bar e_2$; hence all $\rho(eye)$ have a common $1$-eigenvector $v_0$. Taking $y=\eps$, $\rho(e)^2v_0=v_0$, and since $\bar\alpha(e)^2\ne1$ this gives $\rho(e)v_0=v_0$. Then $\rho(e)\rho(y)\rho(e)v_0=v_0$ yields $\rho(y)v_0=v_0$ for all $y\in\cY$, in particular for $y=b_\ell^p$, $\ell\in I$. But the $1$-eigenline of $\rho(b_\ell^p)$ is that of the operator induced by $M_{b_\ell}$, spanned by $\bar e_1+\frac{\beta(\ell)}{1-\alpha(\ell)}\bar e_2$, and by~(iii) these lines are pairwise distinct for infinitely many $\ell\in I$, a contradiction. So there are $z,z'\in\cZ_0$ with $\rho(z)\rho(z')\ne\rho(z')\rho(z)$. All elements of $\rho(\cY)$ are triangular in the basis $(\bar e_2,\bar e_1)$ with diagonal $(\bar\alpha(\cdot),1)$, so the commutator $\rho(z)\rho(z')\rho(z)^{-1}\rho(z')^{-1}$ is a unipotent element of $\cH$, the closure of the group generated by $\rho(\cZ_0)$, and it is not the identity.

\emph{Step 5: inverse weight.} Apply Proposition~\ref{prop:invweight} with $D_0=0$, $D_1=\Q\bar e_2$, $D_2=D$, so that $\kappa_1=\bar\alpha$ and $\kappa_2=1$. There is a free coordinate $i$ of $L$ (a register index) with $\alpha_i(\zeta)\bar\alpha(\zeta)=1$ for all $\zeta\in\cZ_0$.

\emph{Step 6: pumping.} For $\ell\ge0$ let $y_\ell=b_{c+p\ell}\,b_c^{p-1}\in\cY$ and $\zeta_\ell=(ey_\ell e)^{k!}\in\cZ_0$. Then $\bar\alpha(\zeta_\ell)=C\alpha(c+p\ell)^{k!}$ with $C=(\bar\alpha(e)^2\alpha(c)^{p-1})^{k!}\ne0$. Fix $t\in L$. By~\eqref{eq:multiplier}, $\alpha_i(\zeta_\ell)=\varphi(\zeta_\ell)$, where $\varphi(v)=P_{q,v}(t+\epsilon_i)_i-P_{q,v}(t)_i$ is the difference of two functions computed by CCRAs (start in $q$ with valuation $t+\epsilon_i$, respectively $t$, and output register $i$); so $\varphi$ is a rational series (Section~\ref{sec:prelim}). Since $b_{c+p\ell}=x_1w^c(w^p)^\ell x_2$, the word $\zeta_\ell$ has the form $v_0(w^p)^\ell v_1\cdots(w^p)^\ell v_{k!}$ with fixed words $v_j$, so $\ell\mapsto\alpha_i(\zeta_\ell)$ agrees with an exponential polynomial $E$ for large $\ell$. Thus $C\alpha(c+p\ell)^{k!}E(\ell)=1$ for all large $\ell$, where $\ell\mapsto C\alpha(c+p\ell)^{k!}$ is a nonconstant polynomial. This contradicts Lemma~\ref{lem:exppoly}.
\end{proof}

\begin{corollary}\label{cor:reversal}
Over the alphabet $\{a,\#\}$, write $w=u_0\#u_1\#\cdots\#u_k$ with $u_m\in a^*$, and let
\[
g(w)=\sum_{i=0}^k\prod_{j<i}|u_j| ,\qquad h(w)=\sum_{i=0}^k\prod_{j<i}\bigl(|u_j|+1\bigr).
\]
The reversals $g^{\mathrm r}(w)=g(w^{\mathrm{rev}})$ and $h^{\mathrm r}$ are computed by CCRA, but $g$ and $h$ are not. Likewise, for the function $F(0^{m_1}1^{k_1}\cdots0^{m_t}1^{k_t})=(\cdots((m_1+k_1)m_2+k_2)\cdots)m_t+k_t$ of~\cite[Example~52]{MPSfull}, where $m_1,k_t\ge0$ and all other exponents are positive (so the decomposition is unique and $F(\eps)=0$), $F$ is computed by a CCRA and its reversal is not. In particular, the class of functions computed by CCRA over $\Q$ is not closed under reversal, for every alphabet with at least two letters.
\end{corollary}
\begin{proof}
\emph{Machines.} Since $g^{\mathrm r}(u_0\#\cdots\#u_k)=\sum_{i=0}^k\prod_{j>i}|u_j|=:G_k$ satisfies $G_k=G_{k-1}|u_k|+1$ and $G_0=1$, the CCRA with one state, registers $x,A$ initialized to $0$, updates $A:=A+1$ on $a$ and $(x,A):=(xA+1,0)$ on $\#$, and output $xA+1$ computes $g^{\mathrm r}$: after reading $u_0\#\cdots\#u_m'$ with $u_m'$ a prefix of $u_m$, $x=G_{m-1}$ (with $G_{-1}=0$) and $A=|u_m'|$. For $h^{\mathrm r}$ initialize and reset $A$ to $1$ instead. For $F$, a two-state control records whether the last letter was $0$ (initially yes); registers $v=1$, $c=0$; on $0$, $c:=c+1$ after a $0$ and $c:=1$ after a $1$; on $1$, $(v,c):=(vc+1,0)$ after a $0$ and $v:=v+1$ after a $1$; output $vc$ after a $0$ and $v$ after a $1$. All updates are copyless.

\emph{Affine contexts.} On row coordinates $(S,Q,P)$ let
\[
M_a=\begin{pmatrix}1&0&0\\0&1&0\\0&1&1\end{pmatrix},\quad
M_\#=\begin{pmatrix}1&0&0\\1&0&1\\0&0&0\end{pmatrix},\quad
M'_\#=\begin{pmatrix}1&0&0\\1&0&1\\1&0&1\end{pmatrix},\quad
\lambda=(1,0,1),\ \gamma=(1,0,0)^{\mathsf T},
\]
that is, $a:(S,Q,P)\mapsto(S,Q+P,P)$, $\#:(S,Q,P)\mapsto(S+Q,0,Q)$ and $\#':(S,Q,P)\mapsto(S+Q+P,0,Q+P)$. Induction on the number of blocks shows that $(\lambda,M_a,M_\#,\gamma)$ represents $g$ and $(\lambda,M_a,M'_\#,\gamma)$ represents $h$: after $u_0\#\cdots\#u_m'$ the state is $(S,|u_m'|P,P)$ with $S$ the sum of the terms $i\le m$ and $P$ the product of the factors contributed by $u_0,\ldots,u_{m-1}$. Both are minimal, since $\lambda,\lambda M_a,\lambda M_\#$ and $\gamma,M_\#\gamma,M_aM_\#\gamma$ are bases, and likewise $\lambda,\lambda M_a,\lambda M'_\#$ and $\gamma,M'_\#\gamma,M_aM'_\#\gamma$. Take $u=x_1=\eps$, $w=a$, $x_2=\#$, $r_1=(1,0,0)$, $r_2=(0,0,1)$. Then $r_1M_{a^\ell\#}=r_1$ and $r_2M_{a^\ell\#}=\ell r_1+\ell r_2$, while $r_2M_{a^\ell\#'}=(\ell+1)(r_1+r_2)$; and $\lambda=r_1+r_2$. So $g$ has an affine context with $\alpha=\beta=T$ and $h$ with $\alpha=\beta=T+1$; in both cases $\beta/(\alpha-1)$ is nonconstant. For $F^{\mathrm r}(w)=F(w^{\mathrm{rev}})$ use row coordinates $(S,R,Q,P)$ with
\[
1:(S,R,Q,P)\mapsto(S+R+Q,\,R+Q,\,0,\,R+Q),\qquad 0:(S,R,Q,P)\mapsto(S,\,0,\,Q+P,\,P),
\]
$\lambda=(0,1,0,1)$ and output $S+Q$; here $R$ is the weight added by each $1$ of the current run of ones, $P$ the product of the completed runs of zeros, and $Q=P$ times the length of the current run of zeros. This representation of $F^{\mathrm r}$ is minimal of dimension $4$. With $u=\eps$, $x_1=w=0$, $x_2=1$, $r_1=(1,0,0,0)$, $r_2=(0,1,0,1)=\lambda$ we get $r_1M_{0^{\ell+1}1}=r_1$ and $r_2M_{0^{\ell+1}1}=(\ell+1)r_1+(\ell+1)r_2$, an affine context with $\alpha=\beta=T+1$. Theorem~\ref{thm:affinecontext} applies in all three cases. For a larger alphabet $\Sigma\supseteq\{a,\#\}$, extend $g$ by $0$ outside $\{a,\#\}^*$: the reversal of the extension is computed by adding a sink state, and the extension is not CCRA-computable by Lemma~\ref{lem:quotients} with $\Sigma'=\{a,\#\}$.
\end{proof}

Polynomiality of $\alpha$ enters the proof of Theorem~\ref{thm:affinecontext} only in Step~6. The following general form shows that what matters is whether the block scalar is a pure exponential along pumping.

\begin{lemma}\label{lem:units}
Let $E,E'$ be exponential polynomials with $E(\ell)E'(\ell)=1$ for all large $\ell$. Then there is $p'\ge1$ such that for every residue $s$ there are $C_s,\mu_s\in\Qb^\times$ with $E(s+p'\ell)=C_s\mu_s^\ell$ for all large $\ell$.
\end{lemma}
\begin{proof}
Let $\Lambda\subseteq\Qb^\times$ be the finitely generated group generated by the bases of $E$ and $E'$, and $p'$ the order of its torsion subgroup. On the residue class $s+p'\N$, both functions are exponential polynomials in $\ell$ with bases in $\Lambda^{p'}$, which is free abelian, say of rank $m$. By (E1), exponential polynomials with bases in $\Lambda^{p'}$, taken as functions of large $\ell$, form a ring isomorphic to the Laurent polynomial ring $\Qb[T][X_1^{\pm1},\ldots,X_m^{\pm1}]$, whose units are the nonzero constant multiples of Laurent monomials. Hence $E(s+p'\ell)$ is a unit of this ring, that is, of the form $C_s\mu_s^\ell$.
\end{proof}

\begin{theorem}[General form]\label{thm:affinegeneral}
Let $u,x_1,w,x_2$ and $r_1,r_2$ satisfy (i) and (ii) of Definition~\ref{def:affinecontext} for some functions $\alpha,\beta:\N\to\Q$; these agree with exponential polynomials for large $\ell$. Suppose that $\alpha(\ell)\ne0$ for all large $\ell$, and that $|\alpha(\ell)|>1$ for all large $\ell$ or $|\alpha(\ell)|<1$ for all large $\ell$; that $\beta/(\alpha-1)$ is not eventually constant on any infinite arithmetic progression; and that on no infinite arithmetic progression does $\alpha$ agree, for large $\ell$, with a function $C\mu^\ell$, $C,\mu\in\Qb^\times$. Then $f$ is not computed by any CCRA over $\Q$.
\end{theorem}
\begin{proof}
With $v_1,v_2$ as in Step~2, $\alpha(\ell)=r_2M_{x_1}M_w^\ell M_{x_2}v_2$ and $\beta(\ell)=r_2M_{x_1}M_w^\ell M_{x_2}v_1$, which agree with exponential polynomials by (E3). The proof of Theorem~\ref{thm:affinecontext} goes through verbatim up to Step~6, with $c$ chosen so that $\alpha$ has no zeros on $I$, and with $|\cdot|<1$ in place of $|\cdot|>1$ in the second case, the hypothesis on $\beta/(\alpha-1)$ being used in Steps~2 and~4. Step~6 yields an exponential polynomial $E$ with $C\alpha(c+p\ell)^{k!}E(\ell)=1$ for large $\ell$. Let $\Lambda$ be the group generated by the bases of $E$ and of the exponential polynomial $\ell\mapsto\alpha(c+p\ell)$, and $p'$ the order of its torsion subgroup. On a residue class $s+p'\N$ both functions have bases in the free abelian group $\Lambda^{p'}$, so, as in the proof of Lemma~\ref{lem:units}, they are elements of a Laurent polynomial ring over $\Qb[T]$ in which $\alpha^{k!}$ is a unit. A divisor of a unit is a unit, so $\alpha(c+p(s+p'\ell))=C''\mu''^\ell$ for large $\ell$, contrary to the hypothesis.
\end{proof}
For instance, block scalars $\alpha(\ell)=2^\ell-1$ or $\alpha(\ell)=\ell2^\ell$ are excluded just like polynomial ones, whereas $\alpha(\ell)=2^\ell$ is allowed; see Remark~\ref{rem:localcond}(2).

Proposition~\ref{prop:lifting} gives a second, independent obstruction, which applies when the two weights coincide and the inverse-weight condition is void.

\begin{theorem}[Unipotent contexts]\label{thm:unipotentcontext}
Let $u,x_1,w,x_2$ be words, $b_\ell=x_1w^\ell x_2$, and $r_1,r_2\in\Q^{1\times d}$ linearly independent rows with $\lambda M_u\in\Q r_1+\Q r_2$, $\lambda M_u\notin\Q r_1$, such that for all $\ell\ge0$
\[
r_1M_{b_\ell}=\kappa(\ell)\,r_1,\qquad r_2M_{b_\ell}=\kappa(\ell)\bigl(s(\ell)\,r_1+r_2\bigr)
\]
with $\kappa(\ell)\ne0$ for all large $\ell$. If on no infinite arithmetic progression the function $s$ agrees, for large $\ell$, with an exponential polynomial, then $f$ is not computed by any CCRA over $\Q$.
\end{theorem}
\begin{proof}
Suppose a CCRA computes $f$. Step~1 of the proof of Theorem~\ref{thm:affinecontext} applies with $c\ge n$ chosen so that $\kappa(\ell)\ne0$ for $\ell\ge c$, and gives $I=c+p\N$, $u'$, $q$ and the monoid $\cY$ of products of $p$ blocks $b_\ell$, $\ell\in I$. Let $X=\Q r_1+\Q r_2$, invariant under each $M_{b_\ell}$. By hypothesis $s(\ell)\ne0$ for some $\ell\in I$; writing $\lambda M_{u'}=\mu_1r_1+\mu_2r_2$ with $\mu_2\ne0$, the row $\lambda M_{u'}M_{b_\ell}^p$ has $r_1$-coefficient $\kappa(\ell)^p(\mu_1+p\,s(\ell)\mu_2)$ and $r_2$-coefficient $\kappa(\ell)^p\mu_2$, so it is independent of $\lambda M_{u'}$. Hence the rows $\lambda M_{u'}M_y$, $y\in\cY$, span $X$, $K_{u'}=X^\perp$, and $D=W/N\cong\Q^d/X^\perp$ is two-dimensional. With $v_1,v_2$ and $\bar e_1,\bar e_2$ as in Step~2, the operator induced by $M_{b_\ell}$ on $\Q^d/X^\perp$ is $\kappa(\ell)(I+s(\ell)\mathcal N)$, where $\mathcal N\bar e_1=\bar e_2$ and $\mathcal N\bar e_2=0$. These operators lie in the closed commutative group $\mathcal G=\{x(I+y\mathcal N):x\ne0\}$, so for $y\in\cY$ we have $\rho(y)=\bar\kappa(y)(I+S(y)\mathcal N)$ with $\bar\kappa$ multiplicative and $S$ additive, $S(b_{\ell_1}\cdots b_{\ell_p})=\sum_js(\ell_j)$; every $\rho(y)$ is invertible.

Apply Proposition~\ref{prop:compression} with a minimizing $e\ne\eps$, and let $\cZ_1$, $F'_L$ and $\cH'$ be as above. The group $\cH'$ lies in $\mathcal G$, so its unipotent elements commute and its maximal tori are contained in the scalars. Hence in Proposition~\ref{prop:lifting}, $t(z)=\bar\kappa(z)I$ and $\log u(z)=S(z)\mathcal N$, and $S(z)$ is a $\Qb$-linear combination of finitely many rational series in $z$, for $z\in\cZ_1$. Let $y_\ell=b_{c+p\ell}b_c^{p-1}$ and $\zeta_\ell=(ey_\ell e)^{m}$ with $m=|F'_L|$, so that $\zeta_\ell\in\cZ_1$. Then
\[
S(\zeta_\ell)=m\bigl(2S(e)+(p-1)s(c)+s(c+p\ell)\bigr),
\]
and, as in Step~6 of the proof of Theorem~\ref{thm:affinecontext}, the left-hand side agrees with an exponential polynomial for large $\ell$. So does $\ell\mapsto s(c+p\ell)$, contrary to the hypothesis.
\end{proof}

\begin{example}[A palindromic obstruction]\label{ex:symmetric}
Over $\{a,\#\}$ let $E(u_0\#\cdots\#u_k)=\sum_{i=0}^k\prod_{j\ne i}(|u_j|+1)$, the elementary symmetric function of degree $k$ in the numbers $|u_j|+1$. It is invariant under permutations of the blocks, so $E(w^{\mathrm{rev}})=E(w)$. On row coordinates $(S,P,Q_S,Q_P)$ it is represented by $a:(S,P,Q_S,Q_P)\mapsto(S,P,Q_S+S,Q_P+P)$, $\#:(S,P,Q_S,Q_P)\mapsto(Q_S+S+P,\,Q_P+P,\,0,\,0)$, $\lambda=(0,1,0,0)$ and output $Q_S+S+P$; this representation is minimal, and every word matrix containing $\#$ has rank two. With $u=x_1=\eps$, $w=a$, $x_2=\#$, $r_1=(1,0,0,0)$ and $r_2=(0,1,0,0)=\lambda$ we get $r_1M_{a^\ell\#}=(\ell+1)r_1$ and $r_2M_{a^\ell\#}=r_1+(\ell+1)r_2$, so $\kappa(\ell)=\ell+1$ and $s(\ell)=1/(\ell+1)$, which is an exponential polynomial on no arithmetic progression by Lemma~\ref{lem:exppoly}. By Theorem~\ref{thm:unipotentcontext}, $E$ is not CCRA-computable. Here the block products generate a commutative group whose Zariski closure is $\{x(I+y\mathcal N)\}\cong\mathbb G_m\times\mathbb G_a$, condition~(C3) holds for every local group, and the unipotent weight is trivial, so Proposition~\ref{prop:invweight} gives nothing: the obstruction is the translation $s(\ell)=1/(\ell+1)$, which a register would have to accumulate.
\end{example}

The claims about representations and machines in Corollary~\ref{cor:reversal} and Example~\ref{ex:symmetric} were also checked by exact computation (Section~\ref{sec:singular-checks}). We record what the earlier necessary conditions say about $g$ and $h$.

\begin{remark}[What the local condition must look like]\label{rem:localcond}
(1) \emph{Maximal subgroups do not detect the obstruction.} Let $\bar M$ be the Zariski closure in $M_3(\Qb)$ of the matrix monoid of $g$ or of $h$. The word matrices are the powers $M_a^n$ and matrices of rank at most two, and all of them fix the row $r_1$; hence $\bar M\subseteq\{M_a^t:t\in\Qb\}\cup\{X:\operatorname{rank}X\le2\}$ and $r_1X=r_1$ for all $X\in\bar M$. Let $\epsilon\in\bar M$ be idempotent and $H_\epsilon$ the group of units of $\epsilon\bar M\epsilon$, which acts faithfully on $\Qb^{1\times3}\epsilon\ni r_1$ and fixes $r_1$. If $\epsilon=I$, $H_\epsilon=\{M_a^t\}$; if $\operatorname{rank}\epsilon=2$, $H_\epsilon$ embeds in $\{r_1\mapsto r_1,\ r\mapsto br_1+ar\}\cong\mathrm{Aff}_1$; if $\operatorname{rank}\epsilon=1$, $H_\epsilon$ is trivial. So every maximal subgroup of $\bar M$ stabilizes a complete flag through the rational row $r_1$ and has commuting unipotent elements. The groups of Corollary~\ref{cor:local} lie in such an $H_\epsilon$ with $\epsilon$ rational, namely the idempotent $\epsilon$ with the same image and kernel as $M_e$, which exists because $\operatorname{rank}M_e^2=\operatorname{rank}M_e$ and lies in $\bar M$, being the identity of the closure of $\{M_e^n:n\ge1\}$ inside the group of matrices with the same image and kernel as $M_e$ (isomorphic to $\GL(R)$), which is a group by the argument of Lemma~\ref{lem:monoid}; indeed $M_eM_zM_e\in\epsilon\bar M\epsilon$ has the same rank as $\epsilon$; for them the flag $\Q r_1\subset\Q^{1\times3}\epsilon$ is rational, so the conclusion of Corollary~\ref{cor:local} is true for $g$ and $h$. For $h$ the monoid is even more regular: every word matrix containing $\#$ has rank exactly two, so each of them, $M'_\#$ included, lies in a subgroup of $M_3(\Q)$. Since transposition preserves condition~(C3) and maps the closure of the monoid of $f$ to that of $f^{\mathrm r}$, \emph{no characterization of CCRA-computability can be invariant under transposition of the minimal representation}, and any combinatorial condition on resets must distinguish left from right.

(2) \emph{The extra condition is arithmetic.} By Theorem~\ref{thm:affinegeneral}, the obstruction is that the block scalar fails to be a pure exponential along pumping. By contrast, if the block scalar is $2^\ell$, as for $\sum_i\prod_{j<i}2^{|u_j|}$, then the inverse weight $2^{-\ell}$ is realized by the register $x:=x/2$, and that series is computed by the CCRA with registers $T,D$ initialized to $1$, updates $(T,D):=(T/2,2D)$ on $a$ and $T:=T+1$ on $\#$, and output $TD$ (indeed it is invertible). In the reversal $g^{\mathrm r}$ the roles of $r_1$ and $r_2$ are exchanged, the required multiplier is $\kappa_2/\kappa_1=|u_j|$ itself, and it is realized by the counter $A$ in $x:=xA+1$.
\end{remark}

\subsection{One invertible letter and one rank-two reset}
The necessary conditions above are sharp in the following class, which contains $h$, $h^{\mathrm r}$ and $E$ (but not $g$, whose empty block is singular). Let $\Sigma=\{a,\#\}$, let $M_a$ be invertible, let $\operatorname{rank}M_\#=2$, and assume that every word containing $\#$ has rank two. Factor $M_\#=CR$ with $C\in\Q^{d\times2}$, $R\in\Q^{2\times d}$ and put $T(n)=RM_a^nC$. Since $M_\#M_a^nM_\#=CT(n)R$, the rank assumption says $T(n)\in\GL_2(\Q)$ for all $n\ge0$. For $y\in\Q^{1\times2}$ we have $(yR)M_{a^n\#}=yT(n)R$, and for $v=a^{n_1}\#\cdots a^{n_m}\#$ we write $T(v)=T(n_1)\cdots T(n_m)$. Let $\mathcal T\subseteq\GL_2(\Qb)$ be the Zariski closure of the group generated by $\{T(n):n\ge0\}$. Changing the factorization conjugates all $T(n)$ by one rational matrix.

\begin{theorem}\label{thm:onereset}
In this class, $f$ is CCRA-computable if and only if the following hold.
\begin{enumerate}
\item[(a)] Every eigenvalue of $M_a$ is a root of a rational number.
\item[(b)] $\mathcal T^\circ$ is triangularizable over $\Q$.
\item[(c)] If (b) holds and $\mathcal T^\circ$ contains a unipotent element $\ne1$, then in a rational basis $(y_1,y_2)$ of $\Q^{1\times2}$ such that $\Q y_2$ is $\mathcal T$-invariant, writing $T(n)=\left(\begin{smallmatrix}a(n)&b(n)\\0&d(n)\end{smallmatrix}\right)$, the sequences $d(n)/a(n)$ and $b(n)/a(n)$ are linear recurrence sequences.
\end{enumerate}
\end{theorem}
Condition~(c) does not depend on the choices: rescaling $y_1,y_2$ multiplies $b/a$ by a constant, and replacing $y_1$ by $y_1+cy_2$ changes $b$ by $c(d-a)$, hence adds $c(d/a-1)$ to $b/a$.
For $h$ (with $y_1=r_2$, $y_2=r_1$) one gets $a=b=n+1$, $d=1$, and $d/a=1/(n+1)$ is not a linear recurrence sequence. For $E$ one gets $a=d=n+1$, $b=1$, and $b/a$ fails. For $h^{\mathrm r}$ one gets $a=1$, $d=n+1$ and a polynomial $b$ (for instance $b=1$ for a suitable factorization), so (c) holds. Conditions (a) and (c) can be checked effectively (we only sketch this): (a) is a test on the irreducible factors of the characteristic polynomial of $M_a$, and (c) amounts to divisibility of explicit exponential polynomials with algebraic bases on residue classes. We do not discuss an algorithm for~(b).

\begin{proof}
\emph{Preliminaries.} (P1) If (b) holds and $\mathcal T^\circ$ contains a unipotent element $\ne1$, then $\mathcal T^\circ$ lies in the group of upper triangular matrices for some rational basis, its unipotent radical $U\cong\mathbb G_a$ is normalized by $\mathcal T$ and fixes a unique line, which is defined over $\Q$; so that line is $\mathcal T$-invariant, a basis as in~(c) exists, and every element of $\mathcal T$ is upper triangular in it. Moreover $a,b,d$ are linear in $M_a^n$, hence linear recurrence sequences whose characteristic roots are eigenvalues of $M_a$. (P2) Let $\cG_a$ be the Zariski closure of $\{M_a^n:n\ge0\}$, a group by Lemma~\ref{lem:monoid}, with component group of order $N_a$. For each residue $s$, $\{M_a^n:n\equiv s\bmod N_a\}$ is dense in $M_a^s\cG_a^\circ$, so $\Phi(X)=RXC$ maps $M_a^s\cG_a^\circ$ into the closure of $\{T(n):n\equiv s\bmod N_a\}$; the intersection of the irreducible set $\Phi(M_a^s\cG_a^\circ)$ with $\GL_2$ is irreducible, contains $T(s)$, and lies in $\mathcal T$, hence in one coset of $\mathcal T^\circ$. Thus $n\mapsto T(n)\mathcal T^\circ$ has period $N_a$. (P3) By~\cite{BFLM22}, a linear recurrence sequence over $\Q$ whose characteristic roots are roots of rational numbers is computed by a CCRA over a one-letter alphabet.

\emph{Sufficiency.} Assume (a)--(c). For $w=a^{n_0}\#a^{n_1}\#\cdots\#a^{n_k}$ with $k\ge1$,
\begin{equation}\label{eq:blockform}
f(w)=s(n_0)\,T(n_1)\cdots T(n_{k-1})\,t(n_k),\qquad s(n)=\lambda M_a^nC,\quad t(n)=RM_a^n\gamma,
\end{equation}
and $f(a^n)=\lambda M_a^n\gamma$. All entries of $s,t,T$ are linear in $M_a^n$, so by~(a) and (P3) each is computed by a one-letter CCRA, and so is every product of such entries. The machine runs, inside each block $a^{n}$, fresh copies of one-letter CCRAs for the finitely many block quantities used below, and combines them when it reads $\#$; it also keeps the values computed in the first block, and at the end it uses the copies for $t(n_k)$ and, if $k=0$, for $\lambda M_a^n\gamma$. It remains to evaluate the middle product copylessly.

If $\mathcal T^\circ$ contains no unipotent element $\ne1$, it is a torus, split over $\Q$ by~(b); choose a rational basis in which it is diagonal. Let $F=\mathcal T/\mathcal T^\circ$ and $\varphi(n)=T(n)\mathcal T^\circ$, periodic by (P2), so that the control can track $\varphi$ of the prefix. As in Section~\ref{sec:normalform} choose products $r_g$ of matrices $T(n)$ with $r_g\mathcal T^\circ=g$, $r_1=I$, and put $c(g,n)=r_gT(n)r_{g\varphi(n)}^{-1}\in\mathcal T^\circ(\Q)$, a diagonal matrix whose entries $c_1(g,n),c_2(g,n)$ are linear in $M_a^n$. Then $T(n_1)\cdots T(n_m)=c(g_0,n_1)\cdots c(g_{m-1},n_m)\,r_{g_m}$, and the middle product contributes $\sum_{i=1,2}Y_i\,s(n_0)E_{ii}r_{g}\,t(n_k)$ with product registers $Y_i:=Y_i\,c_i(g,n)$. Expanding the bilinear form, each monomial $Y_is_p(n_0)t_q(n_k)$ uses its own copies, and the machine is copyless.

If $\mathcal T^\circ$ contains a unipotent element $\ne1$, use the basis of~(c). With $A=\prod_ja(n_j)$, $D=\prod_jd(n_j)$ and $B$ the upper right entry of the product over $j=1,\ldots,k-1$, we have $B=AV$, where $V$ satisfies the Horner recursion $V:=V\cdot\frac{d(n)}{a(n)}+\frac{b(n)}{a(n)}$ at each $\#$, starting from $V=0$. Hence
\[
f(w)=s_1(n_0)t_1(n_k)A+s_1(n_0)t_2(n_k)AV+s_2(n_0)t_2(n_k)D .
\]
The sequences $a,d$ have characteristic roots among the eigenvalues of $M_a$, and $q=d/a$ is a linear recurrence sequence by~(c). Here the characteristic roots of a linear recurrence sequence are the bases that occur with nonzero coefficient in its exponential-polynomial form for large $n$. Let $\Gamma\subseteq\Qb^\times$ be the group generated by the eigenvalues of $M_a$ and the characteristic roots of $q$ and of $b/a$, and let $p'$ be a multiple of the order of the torsion subgroup of $\Gamma$, which is finite because $\Gamma$ is a finitely generated abelian group. Fix a residue $s$ and write $n=s+p'm$. As functions of $m$, the sequences $a$ and $d$ are exponential polynomials with bases in the group $\Gamma_0$ generated by the $p'$-th powers of the eigenvalues of $M_a$, and $q$ is one with bases in $\Gamma^{p'}=\{\gamma^{p'}:\gamma\in\Gamma\}\supseteq\Gamma_0$. Both groups are torsion-free: if $(\gamma^{p'})^j=1$ for some $j\ge1$, then $\gamma$ is a torsion element, so $\gamma^{p'}=1$. With $K=\Qb[T]$, (E1) identifies exponential polynomials in $m$ with bases in $\Gamma^{p'}$ with elements of the group ring $K[\Gamma^{p'}]$. This ring is a free $K[\Gamma_0]$-module on the cosets of $\Gamma_0$, and multiplication by $a$, a nonzero element of $K[\Gamma_0]$ (all $a(n)\ne0$ because $T(n)$ is invertible), preserves the coset components. Since $qa=d$ has all its mass on the trivial coset and $K[\Gamma_0]$ is a domain (the group ring of a torsion-free abelian group over a domain), all nontrivial coset components of $q$ vanish. So on every residue class the bases of $q$, as a function of $m$, lie in $\Gamma_0$, which consists of roots of rationals by~(a). A characteristic root $\mu$ of $q$ cannot vanish from all residue classes: the characteristic roots with the same $p'$-th power as $\mu$ are among the $\mu\zeta$ with $\zeta^{p'}=1$, and vanishing on every class would give $\sum_\zeta p_{\mu\zeta}(n)\zeta^n=0$ for all large $n$, hence all $p_{\mu\zeta}=0$ by (E1). So $\mu^{p'}\in\Gamma_0$, and the characteristic roots of $q$ are roots of rationals. The same argument applies to $b/a$ (whose characteristic roots were included in $\Gamma$), and a rational sequence satisfying a recurrence over $\Qb$ satisfies one over $\Q$. So all quantities are computed by one-letter CCRAs, and the displayed formula is evaluated with two copies of $A$, $s_1(n_0)$ and $t_2(n_k)$ and one copy of $V$ and $D$; all updates are copyless.

\emph{Necessity.} Assume that a CCRA computes $f$. Corollary~\ref{cor:local} with $e=\eps$ and $Z=\{a\}$ shows that $\cG_a^\circ$ is triangularizable over $\Q$; as $M_a^{N_a}\in\cG_a^\circ$, its eigenvalues are rational, which gives~(a).

For (b) and (c) we run the construction of Section~\ref{sec:control} on the block words $(a^*\#)^*$. Choose a nonempty block word $t_0$ for which $\delta(S,t_0)$ has minimum size among block words, let $e_0=t_0^{N}$ with $N=|\delta(S,t_0)|!$, $R_0=\delta(S,e_0)$, and let $L_0$ be the monoid of block words fixing $R_0$ pointwise (the word $e_0$ for the control is not to be confused with the element $e$ chosen below by the compression); Lemma~\ref{lem:return} holds for block words. By the argument of Lemma~\ref{lem:dense}, with Lemma~\ref{lem:monoid} applied to the dense monoid $\{T(v)\}$ of $\mathcal T$, the closure $\cK_0$ of $\{T(v):v\in L_0\}$ is a group containing $\mathcal T^\circ$. For every word $u$, $\delta(s_0,ue_0)\in R_0$, and the span $X_u$ of the rows $\lambda M_{ue_0}M_v$, $v\in L_0$, lies in $\Q^{1\times2}R$ (as $e_0$ ends with $\#$) and is $\cK_0$-invariant. Lemma~\ref{lem:local} with $L_0$ in place of $L_{S_{\mathrm{ret}}}$ gives a local quotient $D_u\cong\Q^d/X_u^\perp$, the dual of $X_u$, on which $v\in L_0$ acts invertibly, through $T(v)$.

\emph{Necessity of (b).} By Proposition~\ref{prop:compression} and Corollary~\ref{cor:onestate} applied to $L_0$ and $D_u$, the closure $\cH_u$ of the group generated by the operators induced on $D_u$ has a $\Q$-triangularizable identity component. The restriction morphism maps $\cK_0$ onto $\cH_u$ by (A2), so it maps the connected group $\mathcal T^\circ$ into $\cH_u^\circ$; hence $\mathcal T^\circ$ stabilizes a complete rational flag of $X_u$ (the dual of a flag of $D_u$). The rows $\lambda M_{ue_0}$ span $\Q^{1\times d}M_{e_0}=\Q^{1\times2}R$, which is therefore the sum of finitely many of the subspaces $X_u$, an equivariant quotient of their direct sum; the image of a complete invariant rational flag of the direct sum under this surjection is an invariant rational chain with successive quotients of dimension at most one, which gives a complete $\mathcal T^\circ$-invariant rational flag of $\Q^{1\times2}R\cong\Q^{1\times2}$. This is~(b).

\emph{Necessity of (c).} Now assume (b) and $U\ne1$. If $\dim X_u\le1$ for every $u$, each nonzero $X_u$ would be the line $\Q y_2R$, the only line invariant under $U$, so all rows $\lambda M_{ue_0}$ would lie in it, contradicting $\dim\Q^{1\times d}M_{e_0}=2$. Fix $u$ with $X_u=\Q^{1\times2}R$, so that $D=D_u$ is two-dimensional and $v\in L_0$ acts on it by the matrix $T(v)$ in the basis dual to $(y_1R,y_2R)$; in this basis the column $\bar e_1$ spans an invariant line with weight $\kappa_1=a(v)$ and the quotient has weight $\kappa_2=d(v)$, where $a(v),d(v)$ are the diagonal entries of $T(v)$. Apply Proposition~\ref{prop:compression} to $L_0$ with a minimizing $e\ne\eps$. The power-map argument once more (with $\Phi(h)=(T(e)hT(e))^{k!}$) shows that the closure of the group generated by $\rho(\cZ_0)$ contains $\mathcal T^\circ\supseteq U$, whose nontrivial elements act nontrivially on $D$. Proposition~\ref{prop:invweight} gives a free coordinate $i$ with $\alpha_i(\zeta)=d(\zeta)/a(\zeta)$ on $\cZ_0$, and Proposition~\ref{prop:lifting} (unipotent elements of $\mathcal T$ commute) together with the computation $\left(\begin{smallmatrix}x&c'(w-x)\\0&w\end{smallmatrix}\right)\left(\begin{smallmatrix}1&\sigma\\0&1\end{smallmatrix}\right)=\left(\begin{smallmatrix}a&b\\0&d\end{smallmatrix}\right)\iff x=a,\ w=d,\ \sigma=\frac ba-c'\bigl(\frac da-1\bigr)$, for the maximal torus $T_H\subseteq\{(\begin{smallmatrix}x&c'(w-x)\\0&w\end{smallmatrix})\}$ of $\cH'^\circ$ (every torus of the upper triangular group is of this form for some $c'$), shows that $b(\zeta)/a(\zeta)$ is, on $\cZ_1$, a $\Qb$-linear combination of rational series in $\zeta$. Here $b(\zeta)$ is the upper right entry of $T(\zeta)$.

Now we build words with a single varying block. Let $x_n=e_0a^n\#e_0$; it permutes $R_0$, and the permutations of $R_0$ induced by words $e_0ve_0$ form a finite group, so we may fix block words $w_\sigma$ such that $y_n=x_n\,e_0w_{\sigma}e_0\in L_0$ when $\sigma$ is the permutation induced by $x_n$. The homomorphism $\Psi:\cZ\to F'_L$ defined before Proposition~\ref{prop:lifting} has as image a subsemigroup of a finite group, hence a subgroup; fix $Z_g\in\cZ$ with $\Psi(Z_g)=g^{-1}$. Put $\zeta_n=ey_ne\,Z_{\Psi(ey_ne)}\in\cZ_1$. Along each residue class of $n$ modulo the control period of $a$ (with $n\ge|S|$), the coefficients of $R_{ey_ne}$ are exponential polynomials in $n$, so the condition that $B_{R_{ey_ne}}$ lies in a given component of $\cG'_L$ is the vanishing of finitely many exponential polynomials; by the Skolem--Mahler--Lech theorem~\cite{EPSW03}, $\Psi(ey_ne)$ and the permutation induced by $x_n$ are eventually periodic in $n$. So there is a period $P$ such that for $n$ large in each residue class modulo $P$, $\zeta_n$ is a fixed word with the single factor $a^n\#$ in the middle, and $T(\zeta_n)=A_1T(n)A_2$ with fixed upper triangular $A_1,A_2$. Then
\[
\frac{d}{a}(T(\zeta_n))=\frac da(A_1)\frac da(A_2)\cdot\frac{d(n)}{a(n)},\qquad
\frac ba(T(\zeta_n))=\Bigl(\frac ba(A_1)\frac{d(n)}{a(n)}+\frac{b(n)}{a(n)}\Bigr)\frac da(A_2)+\frac ba(A_2),
\]
by the Horner rule $\frac ba(XY)=\frac ba(X)\frac da(Y)+\frac ba(Y)$ for upper triangular $X,Y$. The left-hand sides are values at $\zeta_n$ of $\Qb$-linear combinations of rational series, hence agree with exponential polynomials in $n$ on each residue class for large $n$. Solving, so do $d(n)/a(n)$ and $b(n)/a(n)$; a rational sequence that agrees with exponential polynomials on all residue classes modulo $P$ for large $n$ is a linear recurrence sequence. This proves~(c).
\end{proof}

\begin{remark}\label{rem:severalletters}
For an alphabet $\Sigma_0\cup\{\#\}$ with several invertible letters and the same rank assumption, $T(z)=RM_zC$ for $z\in\Sigma_0^*$, and the sufficiency proof of Theorem~\ref{thm:onereset} goes through (with (P2) replaced by the same argument for the cosets of the identity component of the closure of the group generated by $\{M_a:a\in\Sigma_0\}$, which a permutation automaton tracks) if condition~(C3) holds for the group generated by $\{M_a:a\in\Sigma_0\}$ and for $\mathcal T$, and, when $\mathcal T^\circ$ has unipotent elements, $z\mapsto d(z)/a(z)$ and $z\mapsto b(z)/a(z)$ are CCRA-computable on $\Sigma_0^*$. The necessity argument uses the Skolem--Mahler--Lech theorem to find words with a single varying block, and we do not know whether it extends.
\end{remark}

\subsection{Registers must interact}
Call a CCRA \emph{affine} if every update has the form $x_i:=\alpha_ix_{\sigma(i)}+\beta_i$ with $\alpha_i,\beta_i\in\Q$ and a permutation $\sigma$ of the registers (both depending on the transition); $\alpha_i=0$, that is, a reset, is allowed, and the outputs may be arbitrary polynomials. By Theorem~\ref{thm:main}, every invertible CCRA series is computed by an affine CCRA.

\begin{proposition}\label{prop:interact}
The CCRA series $g^{\mathrm r}$ of Corollary~\ref{cor:reversal} is not computed by any affine CCRA.
\end{proposition}
\begin{proof}
Suppose an affine CCRA with $n$ states and $r$ registers computes $g^{\mathrm r}$, and let $\Delta$ bound the degrees of its output polynomials. Put $p=n!\,r!$ and $c=n$. For every state $s$, $\delta(s,a^n)$ lies on an $a$-cycle, so the composite update over $a^p$ from it is $x_i\mapsto\mu_ix_i+\nu_i$ (its permutation part is the $r!$-th power of a permutation). Hence the composite update along $a^{c+p\ell}$ from $s$ is $x_i\mapsto A_{s,i}(\ell)x_{\sigma_s(i)}+B_{s,i}(\ell)$ with $\sigma_s$ independent of $\ell$, $A_{s,i}(\ell)=C_{s,i}\mu_{s,i}^\ell$ for some constants $C_{s,i},\mu_{s,i}$ (possibly $0$), and $B_{s,i}$ agreeing with an exponential polynomial for $\ell\ge1$ (a factor $0^\ell$ coming from a reset on the cycle vanishes for $\ell\ge1$); and $\delta(s,a^{c+p\ell})$ does not depend on $\ell$. Below all $\ell_j\ge1$, and (E1) is applied on $\{1,2,\ldots\}^k$.

Fix $k>\Delta$ and consider $w(\ell_1,\ldots,\ell_k)=\#a^{c+p\ell_1}\#a^{c+p\ell_2}\cdots\#a^{c+p\ell_k}$; its control run does not depend on $\ell$. Call a function on $\N^k$ \emph{$m$-sparse} if it is a finite sum of products $\prod_{j=1}^kE_j(\ell_j)$ of exponential polynomials in which at most $m$ factors are not of the form $C\mu^{\ell_j}$. By induction along the word, every register value after reading $w(\ell)$ is $1$-sparse: an affine update $x\mapsto A(\ell_j)x'+B(\ell_j)$ multiplies each summand of a $1$-sparse function by a function $C\mu^{\ell_j}$ of the current block and adds a function of $\ell_j$ alone, and the updates on $\#$ have constant coefficients. Hence $g^{\mathrm r}(w(\ell))$ is $\Delta$-sparse, so by (E1) its expansion in the functions $\prod_j\ell_j^{s_j}\mu_j^{\ell_j}$ only involves functions with $s_j\ne0$ for at most $\Delta$ indices $j$. But $g^{\mathrm r}(w(\ell))=\sum_{i=0}^k\prod_{j=i+1}^k(c+p\ell_j)$ is a polynomial in which the monomial $\ell_1\cdots\ell_k$ has coefficient $p^k\ne0$, a contradiction by (E1).
\end{proof}
So outside the invertible sector the normal form~(C2) of Theorem~\ref{thm:main} fails, even with resets: the product $xA$ of two registers in the machine for $g^{\mathrm r}$ is unavoidable. The same proof applies to $h^{\mathrm r}$ and to $F$.

\subsection{Discussion}\label{sec:singular-checks}
\emph{Compression with resets.} For the machine computing $g^{\mathrm r}$ in Corollary~\ref{cor:reversal}, the word $\#\#$ maps every valuation to $(x,A)=(1,0)$, so a minimizing word in Proposition~\ref{prop:compression} has $k=0$; the operator $\bar E$ then has rank one on the three-dimensional minimal representation, and global compression yields nothing. The correct replacement, used in Corollary~\ref{cor:local} and in the proof of Theorem~\ref{thm:affinecontext}, is to compress only monoids of return words that act invertibly on a local quotient, and to record, through Proposition~\ref{prop:invweight}, which inverse weights the compressed registers must carry.

\emph{A conjectural picture.} Theorems~\ref{thm:main}, \ref{thm:rankone}, \ref{thm:affinegeneral}, \ref{thm:unipotentcontext} and~\ref{thm:onereset} suggest the following shape for a characterization. For every monoid $\cY$ of return words that acts invertibly on a local quotient, the induced representation factors, on the semigroup $e\cY e$ for a word $e$ depending on the machine and after passing to a fixed power of each element, through a morphism of algebraic groups from a group of diagonal affine maps, in such a way that the multipliers and translations of the lifted maps are rational series of the words; this machine-dependent form is what Propositions~\ref{prop:invweight} and~\ref{prop:lifting} prove. This condition implies condition~(C3) for the local groups, the inverse-weight condition and the translation condition. Condition~(C3) for local groups together with trivial inverse weights does not suffice, even for single-rank resets: the function $E$ of Example~\ref{ex:symmetric} satisfies (C3) for all its local groups and has $\kappa_2/\kappa_1=1$, yet is not CCRA-computable. Theorem~\ref{thm:onereset} confirms this picture for one invertible letter and one rank-two reset, where conditions (a)--(c) are the intrinsic form of the lifting condition. We conjecture that an intrinsic version of the full lifting condition is sufficient in general, at least when all words containing a singular letter have the same rank: there the block products are invertible, and the cocycle construction of Section~\ref{sec:normalform} should run at the level of blocks, with the torus coordinates, their required inverses, and the translation coordinates computed inside each block. The lifting condition as stated depends on the machine through $e$; an intrinsic formulation, and a decision procedure, are open.

\emph{Exact checks.} Exact computations (Python, \texttt{fractions} and \texttt{sympy}) verified the machines of Corollary~\ref{cor:reversal} on all words of length at most~$12$, and verified symbolically the representations, minimality claims and context identities used above. They also found affine contexts in the reversals of randomly generated CCRAs. For example, for the one-state CCRA over $\{a,b\}$ with registers $(x_0,x_1)$ initialized to $(1,1)$, updates $(x_0,x_1):=(x_1-1,x_0)$ on $a$ and $(x_0,x_1):=(x_0x_1-1,\frac12)$ on $b$, and output $\frac12x_1+1$, the minimal representation of the reversal has dimension $4$ and an affine context with $u=ab$, $x_1=\eps$, $w=aa$, $x_2=b$, $\alpha(\ell)=\frac12-\ell$, $\beta$ quadratic and $\beta/(\alpha-1)$ nonconstant (checked symbolically), so non-closure under reversal is not confined to hand-made examples. These checks support, but are not part of, the proofs.
